% !TeX root = ../papers/nc-p3.tex
\section{十五次元成分の一般中心の孤立特異性と四次K3切断}
\label{sec:finite-central-k3-general-section}
\begin{theorem}\label{thm:finite-central-k3-general-section}
十五次元混合成分のある稠密開集合では、射影中心の一般超平面切断は滑らかな四次K3曲面であり、
射影中心の特異点は有限である。一般窓の既知の二十節点以外の、次元ベクトル $(2,2,2,2)$
を持つ単純も、次数による倍率を除いて高々有限個である。
一般超平面上の次数4のAzumaya隅、四つのねじれ階数4束と偏極から、
元の窓の四層切断を全ての積つきで取り戻せる。
\end{theorem}
\begin{proof}
まず十四次元片側対称枝の一般の窓に一般の二平面を印す。
この十八次元の族では、幾何学的一般の $(T,h,F)$ のPicard格子は $h,F$ で生成される。
実際 $(T,h,F)$ から、一般の繊維でインスタントン束を復元し、鉛筆から切断二平面を復元できる。
一繊維を零点とする切断の直線はSerre拡大の一意性と束の単純性により一意である。
よって偏極K3のモジュライでの像次元は18である。一つの偏極K3上のこの数値の繊維類は
Hodge指数定理で高々有限個である。複素数上のK3の局所周期写像で三独立な因子を保つ変形は
次元17以下なので、可算個のそうした軌跡は幾何学的一般点を含まない。
これは独立な三つの第一Chern類が二形式の周期に独立な三線型条件を課すことによる。
基礎体の変更は幾何学的一般点で行える。生成された格子の判別式の絶対値は64で、
真の飽和があれば指数に2を含む。$F/2$ は繊維の原始性で存在せず、
$(h+jF)/2$ の平方は奇数でK3の偶格子に入らない。これで格子そのものも確定する。

対称枝のこの幾何学的一般点は定理\ref{thm:symmetric-center-cover}の開近傍にも属する。
ここでは四次中心が五次元で、中心のHilbert級数は換算次数で
$(1+t+t^2+t^3)/(1-t)^4$ である。隅の商は可除で、合成可能な非零元の積は零でない。
定理\ref{thm:finite-minimal-picard-regularity}の構成を使い、同所の大切断の四射影に
対応するねじれ束を $V_i$ とする。共通移動を除き、
$\Theta=S^{-1}[3]$ の作用は制限関手 $J$ の下で $L=a+b$ による捩りとなる。
大切断の射を反復して得る軌道代数の各隅を、順序を合わせて
$\mathrm{Hom}(p_i,\Theta^m p_j)$ と書く。これは元の折り畳み経路代数である。
実際定理\ref{thm:cut-serre}と命題\ref{prop:two-window-ext-complex}により
戻り辺一回の両側加群の導来テンソル冪は純粋で通常のテンソル冪と一致し、
合成もその連接である。

随伴三角の連結射 $S[-3]\to\mathrm{Id}$ を $\Theta p_i$ で評価すると
長さ4の四ループ $g_i:p_i\to\Theta p_i$ を得る。
その和を $g$ とする。$J$ と $J\Theta=LJ$ で誘導される評価は乗法と両立し、
各隅でその像は $\mathrm{Hom}(V_i,V_j(mL))$、核は $g$ を始点に付けた短い経路である。
実際 $m\ge1$ は前節の短完全となった随伴三角であり、負次数の束の間のExtはない。
同じ事実は $m=0$ の次数零の射にも成立する。単射は $m=1,i=j$ で $g_i\ne0$ も示す。
連結射の自然性により次数零の大切断の全矢と $g$ は可換である。
また終点に $g$ を付けた像も評価核に入る。この端からの積も非零経路の積の性質で単射で、
次元が等しいので $gB=Bg$。従って $g$ は正則正規元である。
対応自己同型は既に三辺の文字を固定する。四番目は三次重なりの双対であり、
四次socleテンソルの直線を保つので、そこでは一つのスカラー倍である。
正規自己同型は $g$ 自身を固定し、各 $g_i$ は戻り辺をちょうど一つ含むので、このスカラーも1である。
従って $g$ は中心元である。

零番隅の四Veronese商はよって $\bigoplus_m H^0(M,\mathcal End(V_0)(mL))$ である。
その中心は $\bigoplus_m H^0(M,mL)$ である。高い捩りの自己準同型束は全点で生成されるため、
中心元はスカラーであり、スカラー係数は局所行列表示から全点で正則である。
元の全中心 $Z$ の零番隅への写像は単射である（一般商で既に単射）。
さらに $Z/gZ$ もこの商隅へ単射である。実際 $z$ の零番成分が $g$ 倍なら、
評価した中心元の他の成分も、零番から各頂点への非零の束写像との可換性で消える。
各成分は中心のため一般点でスカラー（高い捩りの各隅の全射を使う）であり、
非零の隅間写像は一般点で非零なので同じスカラーを持つからである。
全体で $z=gb$ と書ければ $g$ の中心性と正則性から $b$ は中心である。
K3上で $L^2=4$、$L$ は豊富なので $h^0(mL)=2+2m^2$ ($m>0$)。
これは先の $Z$ のHilbert級数を $(1-t)$ 倍した全係数と一致する。
従って $\operatorname{Proj}(Z/gZ)=M$ は滑らかな曲面である。

中心は定理\ref{thm:symmetric-center-cover}の近傍で重み付き二次・四次完全交叉の平坦族に入り、
四次中心は階数5の束である。超平面切断が滑らかな曲面となる条件は固有性と滑らかさの開性で開である。
上で一つの幾何点と超平面が得られたので、十五次元成分の一般点とその一般超平面で成立する。
対称因子の外の一般点では中心は普通の四次超曲面である。もし射影特異集合が正次元なら、
超平面と会い、その点で切断の接次元は少なくとも3なので滑らかでない。従って特異集合は有限である。
各特異指標のファイバーは有限次元代数で単純は有限個である。
最後に滑らかな超平面を割った隅は、残る三中心変数上で自由で深さ3を持つ。
従って次数別に全切断加群であり、第一コホモロジーも全捩りで消える。
Azumaya層のねじれた分裂束から四冪等元の部分 $V_i$ を取れば、
$n<4$ の積は中心の四次式を割っても変わらず、元の大切断を復元する。
\end{proof}

\begin{proposition}\label{prop:extra-central-singular-characters}
一般の十五次元成分には、既知の二十節点のほかに特異中心指標が実在する。
特に次元ベクトル $(2,2,2,2)$ の単純が存在する。
全ての残りの特異点も通常二重点で、それぞれ非同型単純を二種類持つ場合には、
その個数はちょうど10である。
\end{proposition}
\begin{proof}
複素数上で計算する。使う条件と一般の結論は有限型のデータから標数零の幾何学的一般点に移る。
射影中心 $X$ 上で、四隅の全てを前定理の次数規約で層化した代数を $\mathcal B$ とする。
その完全圏の周期巡回ホモロジーの偶次元数から奇次元数を引いた値は4である。
理由を詳しく述べる。$\operatorname{qgr} B$ は $\mathcal B$ の連接加群圏に同値である。
実際、中心の四次Veronese次数で隅をまとめた有限次数加群の層化と、
高い捩りの全切断を取る逆操作は、通常の有限生成中心に対するSerre対応そのものである。
四つの頂点成分と固定した経路剰余次数を取れば、ここでは整数頂点の一つの連結な圏を使っている。
四つの連続射影の像には高次の順方向Extがなく、逆方向のExtは全て零である。
実際それらの間の層Homは対応する隅の次数付き極大Cohen--Macaulay加群であり、
局所的な射影性と深さ4から中間コホモロジーが消える。
最高コホモロジーは中心のGorenstein双対で計算できる。
射影四次超曲面の双対化層は $\mathcal O_X(-1)$、また命題\ref{prop:smooth-azumaya}のtraceで
隅の双対は反対隅になる。従って最高次の双対は逆順でさらに経路長4を引いた切断であり、
四つの連続頂点の間では常に負の長さで零となる。
次数零の順方向射は元の経路、対角は定数のみである。
四つの単純の線型Koszul分解を層化すると両端が一つずつの射影の完全列になり、
四連続の射影は両向きの全移動を生成する。有限生成加群は有限自由射影分解を持つため
四対象は充満例外列をなす。そのため局所化・加法性から周期巡回のEuler数は4となる。

もし特異点が標準の節点模型
$\operatorname{End}_R(R^a\oplus I^b)$ ($a,b>0$) を持つ $N$ 点のみならば、
同じEuler数は $-56+2N$ である。この計算は大域的な小解消の射影性に依存しない。
滑らかな開集合上はAzumayaなので周期巡回ホモロジーは通常のde Rhamコホモロジーの偶奇化と一致する
（行列traceのエタール降下と滑らかな場合のHochschild--Kostant--Rosenberg同型）。
閉台での局所化列により、これに各点の台つき圏のEuler数を加えればよい。
台つき有界圏は完備化で変わらない。各コホモロジーは有限長なのでこれは有限長加群への平坦基底変換と
有限項の射影解消によるExt比較から従う。局所大域次元は有限である
（元の次数付き正則環を可逆中心元で局所化して次数方向を除き、さらに完備化する）。
従って台つき完全圏を計算している。

標準節点の台つきEuler数は2である。通常のアフィン節点の小解消
$Y=\operatorname{Tot}_{\mathbb P^1}(\mathcal O(-1)^{\oplus2})$ 上で
$\mathcal O\oplus\mathcal O(1)$ の自己準同型環は $\operatorname{End}(R\oplus I)$ のアフィン版である。
実際零切断の外では二冪等元と階数一直線束の表示で一致し、
大域切断はその余次元2の補集合から延長する。高次自己Extは束射影で
$\mathcal O(n+j)$、$n\ge0,\ j=0,\pm1$ の第一コホモロジーに帰着して零で、
基底の例外対の引戻しはアフィン射ゆえ生成する。従ってこのtilting同値で節点台は
$\mathbb P^1$ 台となる。滑らかな $Y$ の開閉局所化によりそのEuler数は
$\chi(Y)-\chi(Y\setminus\mathbb P^1)=2$。
用いている加法性はdg完全圏の局所化に対する混合複体の完全三角を周期化したものであり、
ここでの各偶奇群は有限次元である。非固有な滑らかな複素多様体でもde RhamのEuler数は
Poincaré双対でコンパクト台Euler数と等しいため開閉差が使える。

滑らかな四次三次元多様体のEuler数は $-56$ である。
接束の全Chern類 $(1+h)^5/(1+4h)$ の三次項を $\int h^3=4$ で積分すればよい。
通常二重点のMilnor繊維は三次元球面と同じホモトピー型でEuler数が零なので、
$N$ 節点の補集合のEuler数も $-56$ である。
これで $4=-56+2N$、ゆえに $N=30$ となる。
既知の二十点は全て上の標準模型を持つため、それだけでは矛盾する。
追加点での単純次元は定理\ref{thm:twenty-central-nodes}の直後の系で決まる。
追加点が節点かつ二単純ならその系で非自由な階数一極大Cohen--Macaulay加群 $I$ の模型を持つ。
その最小正方表示は行列式の重複度2より二行二列で、節点の非退化二次主要部により
四成分の一次部は独立である。従ってこれは標準節点模型で、上の計数を適用できる。
\end{proof}
